\chapter{Thank You}\label{cha:thanks}\markboth{Thank You}{Thank You}
% Every entry starts at the margin: a name half-indented reads as an
% accident, and in a credits list the eye wants one left edge.
\begingroup
\setlength{\parindent}{0pt}
\setlength{\parskip}{5pt plus 1.5pt minus 1pt}

The BMC-K4510 is a machine that never existed. Almost everything in it
that \emph{does} exist was written by somebody else first, and given
away. This chapter is the thanks; the Licences chapter (page~\pageref{cha:licences})
is the paperwork.

\section{The heart of the machine}
\textbf{Gábor Lénárt} (LGB) wrote the 45GS02 CPU core in
\href{https://github.com/lgblgblgb/xemu}{Xemu}, the MEGA65 emulator.
It runs here unchanged --- not a line altered, not a bug fixed --- and
every instruction this book describes is his code executing. There
would be no K4510 without it.

\textbf{Dag Lem} wrote reSID, the SID emulation that has been the
reference for thirty years. The K4510's four sound chips are four
copies of it (the 6581 model, as shipped in VICE~3.3).

\textbf{The VICE team}, whose decades of emulation scholarship ---
documentation, test programs, arguments settled in code --- this
project consulted at nearly every turn.

\section{The tongues}
\textbf{Lee Davison} (1966--2013) wrote EhBASIC, the machine's first
BASIC and still the one that comes up in ROM. He is not here to be
asked about the graphics and sound words bolted onto it; the machine
carries his name in its \verb|README| and its startup banner instead.

\textbf{Richard T. Russell} wrote BBC BASIC, and has kept writing it
for forty years. The console edition (BBCTTY) is what runs on the
Tube co-processor in Chapter~\ref{cha:tube}. The name ``BBC BASIC'' is
used here \emph{by his permission}, which is a courtesy this project
does not take lightly.

\textbf{Marcelo Dantas} (``Mockba the Borg'') wrote
\href{https://github.com/MockbaTheBorg/RunCPM}{RunCPM}, which is the
whole of Chapter~\ref{cha:cpm}: a complete CP/M~2.2 with its own CCP,
vendored here unmodified and simply handed a Z80's worth of address
space.

\textbf{Scot W. Stevenson}, \textbf{Sam Colwell} and \textbf{Patrick
Surry} wrote \href{https://github.com/SamCoVT/TaliForth2}{Tali
Forth~2} and put it in the public domain --- a Forth written to be
read, which is why Chapter~\ref{cha:forth} can be honest about how it
works.

\textbf{Tomasz Biela} (``tebe'') wrote
\href{https://github.com/tebe6502/Mad-Pascal}{Mad Pascal} and
\href{https://github.com/tebe6502/Mad-Assembler}{MADS}; the K4510
target in \verb|pascal/| is a guest in his compiler.
\textbf{Wojciech Bociański} (``bocianu''), whose Neo6502 target showed
how such a guest should behave.

\textbf{Steve Wozniak} wrote Wozmon in 1976 in 256 bytes, and the
monitor in this machine is still recognisably it.

\textbf{Michael Steil} maintains
\href{https://github.com/mist64/msbasic}{msbasic}, and
\textbf{Microsoft} released 6502 BASIC~1.1 under the MIT licence in
2025 --- between them, the machine's next native BASIC.

\section{Letters on the screen}
Fonts are the part of a computer you look at longest, and clean ones
with a clear pedigree are hard to come by.

\textbf{The Linux kernel} contributors --- the 8x8 console font
(\verb|font_8x8.c|) that got the text mode on its feet and is still the
default.

\textbf{Paul Gardner-Stephen} and \textbf{Roman Standzikowski}
(FeralChild64) --- the clean-room C64/C65 chargen from
\href{https://github.com/MEGA65/open-roms}{MEGA65 open-roms}, and
\textbf{Retrofan}, whose PXLfont travels with it by permission.

\textbf{Ville-Matias Heikkilä} (``Viznut'') --- unscii, placed in the
public domain.

\textbf{Damian Vila} --- \href{https://codeberg.org/Dmian/font-bescii}{BESCII},
the PETSCII spirit with a clean pedigree, CC0.

\section{The bare metal}
\textbf{Randy Rossi} wrote
\href{https://github.com/randyrossi/bmc64}{BMC64}: VICE on a Raspberry
Pi with no operating system under it, 50\,Hz smooth scrolling and input
latency measured in single frames. It is the direct reason this machine
has a Pi port at all --- the route to the Pi was always meant to be
BMC64's \verb|emux_api| seam, and the bare-metal case it proved is what
made the port thinkable. He also built the \textbf{VIC-II Kawari}, a
modern drop-in VIC-II with video modes the original never had: the
demonstration that you may extend an 8-bit machine's video chip and
still be working inside the tradition rather than outside it. VICKY is
a more reckless cousin of that idea. \textbf{minch} (aminch) has taken
BMC64 over from him and carries it onto the newer Pis; the project is
in good hands.

\textbf{Rene Stange} wrote \href{https://github.com/rsta2/circle}{Circle},
the bare-metal C++ environment that is the entire reason a Raspberry
Pi~3B+ can boot straight into this machine with no operating system
under it. \textbf{Xalior} wrote
\href{https://github.com/Xalior/circle-libsdl2}{circle-libsdl2}, which
let the desktop emulator move to the Pi essentially as written.

\section{The workshop}
The machine is built with \textbf{cc65} (compiler, assembler, linker
and runtime for the whole system ROM), \textbf{64tass}, \textbf{NASM},
\textbf{GCC} and \textbf{GNU Make}; it shows itself to you through
\textbf{SDL2}. This book is set in \textbf{XeLaTeX} with
\textbf{Clear Sans} and \textbf{Iosevka}, and every screenshot in it
was captured from the machine actually running --- a picture that
cannot be produced fails the build.

\section{Heritage}
Some debts are not code.

\textbf{Acorn Computers} --- the Tube, the sideways-ROM model, the
\verb|*| command prefix, and the idea that a second processor should
be a normal thing to own. The Beeb's ghost is all over this machine.

\textbf{Commodore} --- the other half of its soul: PETSCII, the SID,
and the line from the PET to the C65 that stopped one machine short of
this one.

\textbf{The MEGA65 project} --- for keeping the 45GS02 and the C65
dream alive in real silicon, and for open-roms.

\textbf{Kim Lemon} and the Lemoners --- \href{https://www.lemon64.com}{Lemon64}
has been the C64's front door since 1998: the games database, the
reviews and scans, the music, and a forum that actually answers.
Twenty-eight years of one person's dedication to a machine's community,
and a good half of the small facts this project needed came from there.

\textbf{The Neo6502 handbook} and the \textbf{MEGA65 User's Guide} ---
the two books this one is modelled on, down to the page size.

\textbf{The High Voltage SID Collection} and the composers in it, whose
tunes are what \verb|SIDPLAY| plays on a good evening. None of that
music is distributed with this machine (see
the Licences chapter); it is theirs.

\begin{aside}
\textbf{And the machine's other author.} Most of the K4510's own code
--- VICKY, SHEILA, the DMA engine, the ROM, K/OS, the Tube ULA, the
editors, the Pi port --- was written in conversation with Claude
(Anthropic's Claude Code), over several months of long sessions with a
build log to prove it. The design decisions are the author's; a great
deal of the typing was not.
\end{aside}

\section{And whoever is missing}
This list was assembled by hand and is certainly incomplete. If your
work is in this machine and your name is not on this page, that is an
error and not a judgement --- open an issue at
\url{https://github.com/mlongval/bmc-k4510} and it will be fixed in the
next printing.
\endgroup
